% Hand-edited after the approved adversarial review and restoration of the
% 70M sparse-execution control, 21 September 2026.
% Mechanisms: Analysis024 observations/041-kernel-appendix.md;
% Run037 observations/001-operation-latency.md; Run039 observations/001-am-port-verdict.md.
% 70M: Analysis028 observations/001-matched-70m-grid.md; 01_collect.py,
% 03_tables.py and 04_controls.py; Run045 results/matched-grid.json.
% 14M T2 control: Run048 observations/001-causal-hz-latency.md;
% 07_reduce.py and 11_engineering_comparison.py.
% Revision: Analysis034/observations/005-kernel-validation-appendix.md.
\section{Kernel implementation and validation}
\label{app:kernel-results}
\label{app}
\label{app:kernel-diagnostics}

\subsection{Sparse execution mechanisms}
The specialized implementations preserve trained weights and thresholding while exploiting zero activations. At $h,z$, empty activation tiles can be identified before loading the corresponding weights, avoiding both weight reads and matrix instructions; sufficiently sparse rows may instead use scalar arithmetic. In the paths summarized in Table~\ref{tab:kernel-paths}, zero detection at $a,m,q,k,v$ follows operand loading, so skipping avoids arithmetic but retains those reads.

\begin{table}[!htb]
\centering\small
\caption{\textbf{Sparse execution rules in the 14M and initial 70M implementations.}
MMA denotes a matrix multiply--accumulate instruction. Tile dimensions are
tokens by input features. All zero checks execute during the forward pass.}
\label{tab:kernel-paths}
\begin{tabularx}{\linewidth}{@{}>{\raggedright\arraybackslash}p{.24\linewidth}>{\raggedright\arraybackslash}X>{\raggedright\arraybackslash}X@{}}
\toprule
Operation & Execution rule & Work avoided \\
\midrule
QKV, FFN-up ($a,m$) & Check activation fragments after both operands are loaded.
& Skip matrix instructions for empty fragments; retain operand reads. \\
FFN-down ($h$), attention-output ($z$) & Check $8\times16$ tiles before loading
weights. Rows with at most two nonzeros use scalar arithmetic when numerical
guards allow; other rows use the dense matrix path.
& Empty tiles avoid weight reads and matrix instructions. Scalar rows read
only weights paired with nonzero activations. \\
Attention ($QK,PV$) & Check operand fragments after loading them.
& Skip matrix instructions if either fragment is zero; retain operand reads,
masking and softmax. \\
\bottomrule
\end{tabularx}
\end{table}

In Table~\ref{tab:kernel-mechanisms}, \emph{bypass} is omitted MMA instructions divided by issued plus omitted instructions, pooled across layers and all 338 validation sequences. It includes padded work and, at $h,z$, instructions replaced by scalar arithmetic. Bypass differs from $\Smodel$, the model-wide fraction of zero-operand scalar products, and is not the fraction of latency saved.

% Source: Analysis038/observations/003-manuscript-integration.md; 04_manuscript_evidence.py.
The initial 70M port of the 14M kernel takes 3.325\,ms on Base versus
1.611\,ms natively. A subsequent optimization pass improves sparse $h,z$
execution and dense components, reducing Base latency to 2.732\,ms while
leaving a 1.121\,ms gap. The main figure retains this partially optimized
measurement series; it does not substitute the separate-session controls below.

The 31M implementation ports that optimized 70M policy (Run045,
\texttt{opt073}) to hidden width 256, FFN width 1,024 and head dimension 32.
It uses paired LayerNorm, native dense input projections, $h,z$ output
projections with empty-tile skipping, scalar execution for rows containing at most eight nonzeros
and a matrix-instruction fallback, dense FlashAttention, and a dense
full-vocabulary output head. Dimension-specific tiling and masks are adapted
without a new optimization search. Base and $T_2/P_h$ use
\texttt{opt073-31m-v1} (Run054); the gate-aware
\texttt{opt073-31m-t7-v1} (Run055) retains the same policy while applying
the seven-site thresholds, including $a,m$ gates in the paired normalization.
Table~\ref{tab:kernel-paths} describes the historical implementations, not
every path in these later variants.

\subsection{Numerical validation and timing protocol}
\label{app:timing-protocol}
Measurements use an RTX\,5090, BF16, batch one, 2,048-token causal sequences,
and full vocabulary logits for every token (PyTorch 2.11.0, Transformers 5.12.1,
CUDA 12.8). This is full-sequence inference with fixed weights and shapes;
cached decoding and concurrent serving are outside the evaluated scope.
Final checkpoints were checked on all 338 validation sequences in three
independent processes. Relative to the unmodified eager reference, acceptance
required finite outputs, $|y-\hat y|\leq0.25+0.02|\hat y|$ for every logit,
relative $\ell_2$ error at most $0.02$ per sequence, and absolute mean
validation-loss difference at most $0.001$. Exact agreement in the 14M control
below is stronger than these general tolerances. At 31M, all eleven final checkpoints pass these criteria on all 338 sequences in three processes each: 18 Base/$T_2/P_h$ processes in Run054 and 15 $T_7/P_h$ processes in Run055 (33 total). This qualification applies to the measured checkpoints, not every possible weight or activation pattern.

PyTorch/SDPA and specialized implementations use CUDA graphs with identical
inputs. We report geometric means of synchronized host timings over 64 fixed
validation sequences, seven randomized paired passes and three independent
processes. Timing includes embeddings, transformer blocks, final normalization,
vocabulary logits, thresholding and all input-dependent sparsity checks.
Compilation, static weight preparation, graph capture and equal input staging
are excluded. Activations are recomputed on every forward pass.

Cross-recipe speedups are $t_{\mathrm{PyTorch},T_0/P_0}/t_{\mathrm{specialized},r}$, using a fixed native Base reference per size. Table~\ref{tab:endpoints-14m-70m} uses 0.65544\,ms at 14M and 1.61105\,ms at 70M; Table~\ref{tab:31m-results} uses 1.02386\,ms at 31M. The specialized 14M Base takes 0.65157\,ms (0.652\,ms rounded), a different denominator from native Base. The 31M Base/$T_2/P_h$ and $T_7/P_h$ timings come from separate sessions on different RTX\,5090 devices, with the Run054 native Base held fixed for both; small cross-session differences should not be overinterpreted. The controls below compare execution configurations of the same checkpoint, with explicit denominators.

\subsection{Controlled 14M execution ablations}
\label{app:t2-hz-control}
Table~\ref{tab:t2-hz-control} isolates sparse $h,z$ execution at $T_2/P_h$, $\kappa=0.1$, holding checkpoint, thresholding, fusion and other paths fixed. Disabling it removes tile skipping and scalar substitution, retaining the fused kernel's dense path. All six layers have identical $h,z$ values and zero masks; logits agree exactly on all 338 sequences in three processes. These full-model controls do not separate memory-read, inspection and arithmetic costs.

\begin{table}[H]
\centering\small
\caption{\textbf{14M $T_2/P_h$ execution control at $\kappa=0.1$.}
All configurations retain thresholding. Full-model geometric-mean latencies
follow the timing protocol above; sparse execution includes skipping and
scalar substitution.}
\label{tab:t2-hz-control}
\begin{tabular}{@{}lccr@{}}
\toprule
Configuration & $h$ sparse execution & $z$ sparse execution & Latency (ms) \\
\midrule
A: full execution & On & On & 0.5588 \\
B: disable $h$ & Off & On & 0.6485 \\
C: disable $z$ & On & Off & 0.5630 \\
D: disable both & Off & Off & 0.6415 \\
\bottomrule
\end{tabular}
\end{table}

Joint savings ($D-A$) are $82.7\,\mu$s, or $12.9\%$ of $D$ ($12.8$--$13.0\%$ across processes). Disabling $h$ adds $89.7\,\mu$s; enabling $z$ saves $4.2\,\mu$s with $h$ on but adds $7.0\,\mu$s with $h$ off, precluding additive attribution. (The unmodified kernel takes 0.5614\,ms, $12.5\%$ below $D$; its $2.6\,\mu$s offset from $A$ limits interpretation of the smaller $z$ effect.)

The same checkpoint takes 0.7018\,ms under PyTorch/SDPA: $A$ reduces latency by $20.4\%$, with $D-A$ accounting for $57.9\%$ of the gap. This includes other implementation differences; the share does not apply to the Base-checkpoint comparison.

% Controls: Analysis028/04_controls.py; data/retained-controls.json.
\subsection{70M execution control}
\label{app:kernel-transfer}
Table~\ref{tab:70m-hz-control} holds the 70M checkpoint, thresholding and other components fixed. Disabling sparse $h,z$ execution removes skipping and scalar substitution, retaining the same fused kernel's dense path. All configurations pass the full-validation criteria above; bitwise agreement was not established.

\begin{table}[H]
\centering\small
\caption{\textbf{70M $T_7/P_h$ execution controls at $\kappa=0.5$.}
Full-model geometric means with weights, thresholding and other optimized
components fixed. The session differs from the main-text latency sweep.}
\label{tab:70m-hz-control}
\begin{tabularx}{\linewidth}{@{}>{\raggedright\arraybackslash}Xr@{}}
\toprule
Execution configuration & Latency (ms) \\
\midrule
Optimized implementation, sparse $h,z$ execution enabled & 1.128 \\
Same implementation, sparse $h,z$ execution disabled & 2.704 \\
Native PyTorch $h,z$ replacement & 1.365 \\
\bottomrule
\end{tabularx}
\end{table}

Sparse $h,z$ execution reduces latency by $58.3\%$ relative to the 2.704\,ms dense path. Against the faster 1.365\,ms native $h,z$ replacement, the reduction is $17.4\%$. This comparison includes fusion and layout effects, not only sparsity exploitation.
